\section{Exact cycle laws for symplectic edge observations}\label{sec:cycles}

Let $n\ge1$ and let $V=\mathbb F_q^{2n}$ carry a nondegenerate alternating bilinear form
$\langle\ ,\ \rangle$, put $N=q^{2n}$, and let $X_v$ be independent uniform
vectors in $V\setminus\{0\}$. An oriented edge $u\to v$ records
$Y_{uv}=\langle X_u,X_v\rangle$. Its law is
\[
 p_N(a)=\begin{cases}(N-q)/(q(N-1)),&a=0,\\
 N/(q(N-1)),&a\ne0.\end{cases}
\]
All coordinates of this law are positive since $N\ge q^2$.
Reversing an edge negates its label. This coordinatewise bijection preserves
$p_N$ and the product reference law, so it preserves total variation and
chi-square divergence. Thus any fixed cycle orientation can be converted
to the cyclic orientation used below.
For a cyclically oriented cycle of length $\ell\ge3$, write $\mu_{\ell,N}$
for the edge-vector law and $\pi_{\ell,N}=p_N^{\otimes\ell}$.
Put $\kappa(0)=q-1$ and $\kappa(a)=-1$ for $a\ne0$.

\subsection{The Fourier law}\label{sec:fourier}
\begin{theorem}[Complete cycle Fourier law]\label{thm:fourier}
Fix any nontrivial additive character $\psi$ of $\mathbb F_q$ and use
$\widehat\mu(t)=\mathbb E\psi(\sum_i t_iY_i)$.
If at least one $t_i$ is zero, then
\[
 \widehat\mu_{\ell,N}(t)=\left(-\frac1{N-1}\right)^{\#\{i:t_i\ne0\}}
 =\widehat\pi_{\ell,N}(t).
\]
If all $t_i\ne0$, then
\[
 \widehat\mu_{2k+1,N}(t)=\frac{N^{k+1}-1}{(N-1)^{2k+1}},
\]
and
\[
 \widehat\mu_{2k,N}(t)=
 \frac{1+N^k\bigl((N-1){\bf1}_{t_1t_3\cdots t_{2k-1}=t_2t_4\cdots t_{2k}}-1\bigr)}
 {(N-1)^{2k}}.
\]

\end{theorem}
\begin{proof}
Forests have independent edge labels: condition on all vertices except a
leaf; its label has the same law $p_N$ for every value of its nonzero
neighbor, and remove leaves inductively. Zero Fourier coordinates remove
edges and leave a forest.
For the full cycle define matrices indexed by $V\setminus\{0\}$,
\[
 F_t(x,y)=\psi(t\langle x,y\rangle),\qquad
 P_c(x,z)={\bf1}_{z=cx},\qquad J(x,z)=1.
\]
For $t\ne0$, additive orthogonality gives $F_t\mathbf1=-\mathbf1$, and
\[
 F_sF_t=NP_{s/t}-J.
\]
The constant subspace and its orthogonal complement $H$ are invariant.
On $H$, pair the factors in the cycle trace. For an even cycle their
product is $N^kP_c$ with
$c=t_1t_3\cdots t_{2k-1}/(t_2t_4\cdots t_{2k})$.
The trace of $P_c$ on $H$ is $(N-1){\bf1}_{c=1}-1$.
For an odd cycle the product is $N^kP_cF_t$ on $H$.
Every diagonal entry of $P_cF_t$ is $1$, because the form is alternating,
so its full trace is $N-1$; its constant eigenvalue is $-1$, and hence
its trace on $H$ is $N$.
Adding the constant eigenvalue $(-1)^\ell$ and dividing by
$(N-1)^\ell$ proves the formulas. No division by $2$ occurs, so the proof
includes characteristic $2$.
\end{proof}

\subsection{Pointwise inversion and the even-cycle resonance}
\begin{theorem}[Exact pointwise cycle law]\label{thm:pmf}
For odd $\ell=2k+1$,
\[
 \mu_{\ell,N}(a)-\pi_{\ell,N}(a)
 =\frac{N^{k+1}}{q^\ell(N-1)^\ell}\prod_{i=1}^\ell\kappa(a_i).
\tag{O}\label{eq:odd-pmf}
\]
For even $\ell=2k$, define the integer-valued function
\[
 K_\ell(a)=
 \sum_{\substack{t_i\in\mathbb F_q^\times\\
 t_1t_3\cdots t_{\ell-1}=t_2t_4\cdots t_\ell}}
 \psi\left(-\sum_i t_ia_i\right).
\]
Then
\[
 \mu_{\ell,N}(a)-\pi_{\ell,N}(a)
 =\frac{N^k}{q^\ell(N-1)^\ell}
 \left((N-1)K_\ell(a)-\prod_i\kappa(a_i)\right).
\tag{E}\label{eq:even-pmf}
\]
The sum defining $K_\ell$ has the explicit real evaluation
\[
 K_\ell(a)=\begin{cases}
 \displaystyle\frac{\prod_i\kappa(a_i)}{q-1},
     &\text{some }a_i=0,\\[6pt]
 \displaystyle\frac{1+q^k((q-1){\bf1}_{\mathcal R(a)}-1)}{q-1},
     &\text{all }a_i\ne0,
 \end{cases}
\]
where
\[
 \mathcal R(a):\quad
 a_1a_3\cdots a_{2k-1}=(-1)^k a_2a_4\cdots a_{2k}.
\]
This is a multiplicative resonance condition on the observed edge labels.

\end{theorem}
\begin{proof}
Additive Fourier inversion gives
\[
 \mu_{\ell,N}(a)-\pi_{\ell,N}(a)
 =q^{-\ell}\sum_{t\in\mathbb F_q^\ell}
 \bigl(\widehat\mu_{\ell,N}(t)-\widehat\pi_{\ell,N}(t)\bigr)
 \psi\left(-\sum_i t_i a_i\right).
\]
By \cref{thm:fourier}, the summand vanishes whenever any $t_i=0$.
For odd $\ell$, the remaining Fourier difference is constant, and
$\sum_{t\ne0}\psi(-ta)=\kappa(a)$ gives (O).
For even $\ell$, separating the balanced-frequency indicator from the
constant term gives (E). It remains to evaluate $K_\ell$.
Insert multiplicative-character orthogonality for the constraint on $t$.
The trivial character contributes $\prod_i\kappa(a_i)/(q-1)$.
If some $a_i=0$, every nontrivial-character term vanishes.
If all $a_i\ne0$, let $G(\chi)=\sum_{x\ne0}\chi(x)\psi(x)$.
The nontrivial-character contribution factors into $k$ copies of each
Gauss sum, with
$G(\chi)G(\chi^{-1})=\chi(-1)q$.
It is therefore
\[
 \frac{q^k}{q-1}\sum_{\chi\ne1}
 \chi\left((-1)^k\frac{a_2a_4\cdots a_{2k}}{a_1a_3\cdots a_{2k-1}}\right).
\]
The last character sum is $q-2$ on $\mathcal R(a)$ and $-1$ elsewhere.
The Gauss identity follows directly by setting $x=uy$ in
$G(\chi)G(\chi^{-1})$ and summing over $y$; thus this inversion needs no
external theorem beyond finite-field character orthogonality.
\end{proof}

\subsection{Exact total variation and chi-square divergences}
\begin{theorem}[Exact one-cycle distances]\label{thm:distances}
Use
$d_{\rm TV}(\mu,\pi)=\frac12\sum_a|\mu(a)-\pi(a)|$ and
$\chi^2(\mu\Vert\pi)=\sum_a(\mu(a)-\pi(a))^2/\pi(a)$.
For odd $\ell=2k+1$,
\[
 d_{\rm TV}(\mu_{\ell,N},\pi_{\ell,N})=
 2^{\ell-1}(1-q^{-1})^\ell\frac{N^{k+1}}{(N-1)^\ell},
\tag{TV-O}
\]
\[
 \chi^2(\mu_{\ell,N}\Vert\pi_{\ell,N})
 =(q-1)^\ell\frac{N}{(N-q)^\ell}.
\tag{CS-O}
\]
For even $\ell=2k$,
\begin{align*}
 d_{\rm TV}(\mu_{\ell,N},\pi_{\ell,N})
 ={}&\frac{N^k(q-1)^{\ell-2}}{q^\ell(N-1)^\ell}\\
 &\quad\times\left[
 (N-q)\bigl((q-1)2^{\ell-1}-(q-2)\bigr)
 +(N-1)(q-2)q^k\right],
\tag{TV-E}
\end{align*}
\[
 \chi^2(\mu_{\ell,N}\Vert\pi_{\ell,N})
 =(q-1)^{\ell-2}
 \left((N-q)^{2-\ell}+(q-2)(N-1)^{2-\ell}\right).
\tag{CS-E}
\]

\end{theorem}
\begin{proof}
\textit{Odd computation.}
Equation (O) and $\sum_a|\kappa(a)|=2(q-1)$ give (TV-O).
For (CS-O), put
\[
 b_N(0)=\frac{(q-1)N}{N-q},\qquad b_N(a)=-1\ (a\ne0).
\]
Then $\mathbb E_{p_N}b_N=0$,
$\mathbb E_{p_N}b_N^2=(q-1)N/(N-q)$, and
\[
 \frac{\mu_{2k+1,N}(a)}{\pi_{2k+1,N}(a)}-1
 =N^{-k}\prod_i b_N(a_i).
\]
Squaring and factoring the expectation proves (CS-O).

\textit{Even computation.}
Write $A=N-q$, $C=(N-1)q^k$ and
$B(a)=(N-1)K_\ell(a)-\prod_i\kappa(a_i)$.
If $z\ge1$ coordinates of $a$ vanish, then
$B(a)=A\prod_i\kappa(a_i)/(q-1)$, so its sign is
$(-1)^z$ (as $\ell$ is even).
If all coordinates are nonzero, then
\[
 B(a)=\begin{cases}(A+(q-2)C)/(q-1),&\mathcal R(a),\\
 (A-C)/(q-1),&\neg\mathcal R(a).
 \end{cases}
\]
The first expression is positive and the second negative whenever its
class is nonempty. There are $(q-1)^{\ell-1}$ resonant vectors and
$(q-2)(q-1)^{\ell-1}$ nonresonant vectors.
For vectors having at least one zero coordinate, the sum of $|B|$ is
\[
 A(q-1)^{\ell-1}(2^\ell-1).
\]
The sum over vectors without zeros is
\[
 (q-1)^{\ell-2}\bigl((3-q)A+2(q-2)C\bigr).
\]
Adding these terms and multiplying by
$N^k/[2q^\ell(N-1)^\ell]$ gives (TV-E).
For chi-square, (E) gives
\[
 \chi^2=\frac1{q^\ell(N-1)^\ell}
 \sum_a\left(\frac N{N-q}\right)^zB(a)^2.
\]
The terms with $z\ge1$ sum to
\[
 (q-1)^{\ell-2}A^2
 \left[\left(\frac{q(N-1)}A\right)^\ell-1\right].
\]
The $z=0$ terms sum to
$(q-1)^{\ell-2}(A^2+(q-2)C^2)$.
The cancellation of the $A^2$ terms and $q^{2k}=q^\ell$ give (CS-E).
\end{proof}

\section{Cactus factorization and independence}\label{sec:cactus}
\begin{theorem}[Exact block independence]\label{thm:cactus}
A cactus graph is a finite simple graph whose distinct cycles have at
most one common vertex. Its blocks containing edges are bridges and simple
cycles; isolated vertices carry no observations.
Under the symplectic model, the edge vectors associated with all these
blocks are mutually independent. Each bridge has law $p_N$, and each
cyclically oriented cycle of length $\ell$ has law $\mu_{\ell,N}$.

\end{theorem}
\begin{proof}
For a block sharing an articulation vertex $v$ with the remainder,
condition on all vectors outside the block except $X_v$.
The remaining random vectors in the block are independent and uniform.
The conditional law of its edge labels given $X_v=x$ is independent of
the chosen nonzero $x$: the symplectic group acts transitively on
$V\setminus\{0\}$ and preserves all labels, while carrying uniform
vectors to uniform vectors. Hence the block labels are independent of
everything in the remaining graph. Remove leaf blocks in the block-cut
tree and iterate. The same argument works componentwise.
\end{proof}

\begin{corollary}[Independence characterizes forests]\label{cor:forest}
For any finite simple graph $G$, its edge labels are independent if and
only if $G$ is a forest.
\end{corollary}
\begin{proof}
Sufficiency was proved above. If $G$ has a cycle,
retain its cycle-edge marginal; it has one of the strictly positive TV
distances (TV-O) or (TV-E), so the complete family cannot be independent.
\end{proof}

\subsection{Paired cycle lengths}
\begin{corollary}[Equal leading information]\label{cor:paired}
For each fixed $q$ and $r\ge1$, both lengths
$\ell=2r+1$ and $\ell=2r+2$ satisfy
\[
 \chi^2(\mu_{\ell,N}\Vert\pi_{\ell,N})
 =(q-1)^{2r+1}N^{-2r}(1+O_{q,r}(N^{-1})).
\]
Moreover, uniformly in $a$,
$\mu_{\ell,N}(a)/\pi_{\ell,N}(a)-1=O_{q,\ell}(N^{-r})$.
\end{corollary}
\begin{proof}
The information asymptotics follow by expanding (CS-O) and (CS-E). In the even case the leading factor is $(q-1)^{2r}(1+q-2)=(q-1)^{2r+1}$.
The uniform likelihood bound follows at once from (O), (E), and $p_N(a)\to1/q$.
\end{proof}

